\documentclass{article}
\usepackage[utf8]{inputenc}
\usepackage{geometry}
\geometry{a4paper, margin=1in}

\title{Respuestas - Evaluación Semana 1 (MX/NS)}
\author{Nathaniel}
\date{\today}

\begin{document}

\maketitle

\section*{Soluciones al Cuestionario}

\begin{enumerate}
    \item \textbf{Función de registros MX y NS:}
    \begin{itemize}
        \item Un registro \textbf{MX} (Mail Exchanger) especifica qué servidor de correo es responsable de recibir los correos electrónicos en nombre de un dominio.
        \item Un registro \textbf{NS} (Name Server) indica cuáles son los servidores de nombres autorizados para una zona DNS, es decir, qué servidores tienen la información oficial sobre ese dominio.
    \end{itemize}

    \item \textbf{Prioridad en registros MX:}
    El número de preferencia o prioridad en un registro MX indica el orden en que los servidores de correo deben ser contactados. El servidor con el número \textbf{más bajo} es el preferido (el principal). Si ese servidor no responde, el cliente intentará conectarse con el servidor que tenga el siguiente número de prioridad más bajo. Esto permite tener servidores de correo de respaldo para garantizar que no se pierdan correos si el servidor principal falla.

    \item \textbf{Mejora de la modularidad de la función:}
    Para hacer la función \texttt{get\_mx\_ns\_records()} más modular, en lugar de imprimir los resultados directamente, debería \textbf{retornar los datos} en una estructura, como una lista de tuplas o un diccionario. Por ejemplo, podría retornar \texttt{\{'MX': [...], 'NS': [...] \}}. De esta forma, la función que la llama (en \texttt{auditorias.py}) es la que se encarga de decidir cómo presentar esa información: imprimirla en la consola, guardarla en un archivo JSON, o mostrarla en un reporte HTML. Esto separa la lógica de obtención de datos de la lógica de presentación.

\end{enumerate}

\end{document}